复共轭与换叶分别作用为
\begin{equation}\label{eq:symmetry-domains}
z\mapsto\conj{z}:D_\pm\longrightarrow D_\mp,\qquad
h:z\mapsto-\dfrac{q_0^2}{z}:D_\pm\longrightarrow D_\mp.
\end{equation}

\paragraph{两种映射怎样作用于区域.}
共轭不改变 $|z|$, 而把 $\operatorname{Im}z$ 变为其相反数,
所以区域定义中的乘积满足
\[
\operatorname{Im}(\conj{z})
\left(1-\dfrac{q_0^2}{|\conj{z}|^2}\right)
=-\operatorname{Im}z\left(1-\dfrac{q_0^2}{|z|^2}\right).
\]
换叶映射则满足
\[
h(z)=-\dfrac{q_0^2\conj{z}}{|z|^2},
\qquad
\operatorname{Im}h(z)=\dfrac{q_0^2}{|z|^2}\operatorname{Im}z,
\qquad
|h(z)|^2=\dfrac{q_0^4}{|z|^2}.
\]
代回同一个区域判别式,
\begin{align*}
\operatorname{Im}h(z)
\left(1-\dfrac{q_0^2}{|h(z)|^2}\right)
&=\dfrac{q_0^2}{|z|^2}\operatorname{Im}z
\left(1-\dfrac{|z|^2}{q_0^2}\right)\\
&=-\operatorname{Im}z\left(1-\dfrac{q_0^2}{|z|^2}\right).
\end{align*}
所以两种映射分别都交换 $D_+$ 与 $D_-$.
它们的对易性可直接计算:
\[
h(\conj{z})=-\dfrac{q_0^2}{\conj{z}}
=\conj{\left(-\dfrac{q_0^2}{z}\right)}=\conj{h(z)}.
\]
连续作用两次换叶还给出
\[
h(h(z))=-\dfrac{q_0^2}{-\dfrac{q_0^2}{z}}=z.
\]
这些关系决定离散谱轨道, 也决定哪些列的等式能够在相同解析域中比较.

二者对易, 其复合保持 $D_\pm$. 以下先在普通连续谱上比较 Jost 归一化, 再将相应列关系延拓到解析域.

\begin{proposition}[Jost 解的共轭与换叶约化]\label{prop:symmetry}
令
\[
\sigma_2=\begin{pmatrix}0&-\ii\\\ii&0\end{pmatrix},\qquad
C(z)=\dfrac{\ii z}{q_0}\sigma_1.
\]
在定理 \ref{thm:jost} 的条件下,
\begin{equation}\label{eq:conjugation}
\phi_\pm(z)=\sigma_2\conj{\phi_\pm(\conj{z})}\sigma_2,\qquad
\mu_\pm(z)=\sigma_2\conj{\mu_\pm(\conj{z})}\sigma_2,
\end{equation}
\begin{equation}\label{eq:musheet}
\phi_\pm(h(z))=\phi_\pm(z)C(z),\qquad
\mu_\pm(h(z))=\mu_\pm(z)C(z).
\end{equation}
逐列写为
\begin{alignat}{2}
&\phi_{\pm,1}(z)=\ii\sigma_2\conj{\phi_{\pm,2}(\conj{z})},&\qquad
&\phi_{\pm,2}(z)=-\ii\sigma_2\conj{\phi_{\pm,1}(\conj{z})},\label{eq:jost-conjugate-columns}\\
&\phi_{\pm,1}(h(z))=\dfrac{\ii z}{q_0}\phi_{\pm,2}(z),&\qquad
&\phi_{\pm,2}(h(z))=\dfrac{\ii z}{q_0}\phi_{\pm,1}(z).\label{eq:jost-sheet-columns}
\end{alignat}
修正列 $\mu_{\pm,j}$ 满足相同关系.
\end{proposition}
\begin{proof}

\par\medskip\noindent\textbf{共轭约化中的矩阵乘法.}\quad 
先计算三个常矩阵的基本关系:
\[
\sigma_2^2
=\begin{pmatrix}(-\ii)\ii&0\\0&\ii(-\ii)\end{pmatrix}=I,
\qquad
\sigma_2\sg
=\begin{pmatrix}0&\ii\\\ii&0\end{pmatrix},
\]
\[
\sigma_2\sg\sigma_2
=\begin{pmatrix}-1&0\\0&1\end{pmatrix}=-\sg,
\qquad
\sigma_2\sigma_1\sigma_2=-\sigma_1.
\]
对势矩阵则有
\begin{align*}
\sigma_2\conj{Q}
&=\begin{pmatrix}0&-\ii\\\ii&0\end{pmatrix}
\begin{pmatrix}0&\conj{q}\\-q&0\end{pmatrix}
=\begin{pmatrix}\ii q&0\\0&\ii\conj{q}\end{pmatrix},\\
\sigma_2\conj{Q}\sigma_2
&=\begin{pmatrix}0&q\\-\conj{q}&0\end{pmatrix}=Q.
\end{align*}
单值化有理函数的系数为实数, 所以
\[
\conj{k(\conj{z})}=k(z),
\qquad
\conj{\lambda(\conj{z})}=\lambda(z),
\qquad
\conj{F(k(\conj{z}))}=F(k(z)).
\]
于是
\begin{align*}
\sigma_2\conj{U(\conj{z})}\sigma_2
&=\sigma_2(\ii k(z)\sg+\conj{Q})\sigma_2\\
&=-\ii k(z)\sg+Q=U(z),\\
\sigma_2\conj{Y(\conj{z})}\sigma_2
&=\sigma_2\left(I+\dfrac{\ii q_0}{z}\sigma_1\right)\sigma_2\\
&=I-\dfrac{\ii q_0}{z}\sigma_1=Y(z).
\end{align*}
相位满足
\[
\conj{\theta(\conj{z})}
=-\conj{\lambda(\conj{z})}
\bigl[x-\conj{F(k(\conj{z}))}t\bigr]
=-\lambda(z)[x-F(k(z))t]=\theta(z).
\]
这里 $x,t$ 为实变量, 这一点在取共轭时也被使用了.

由 $\sigma_2\sg\sigma_2=-\sg$、$\sigma_2\conj{Q}\sigma_2=Q$ 及 $k(\conj{z})=\conj{k(z)}$,
\[
U(z)=\sigma_2\conj{U(\conj{z})}\sigma_2,\qquad
Y(z)=\sigma_2\conj{Y(\conj{z})}\sigma_2.
\]
又有 $\theta(\conj{z})=\conj{\theta(z)}$, 故
\[
\sigma_2\conj{Y(\conj{z})e^{\ii\theta(\conj{z})\sg}}\sigma_2
=Y(z)e^{\ii\theta(z)\sg}.
\]

\par\medskip\noindent\textbf{共轭指数的负号与列交换的负号.}\quad 
取共轭先将 $\ii$ 变为 $-\ii$, 故
\[
\conj{e^{\ii\theta(\conj{z})\sg}}
=e^{-\ii\theta(z)\sg}.
\]
再与 $\sigma_2$ 共轭, 对角元交换:
\begin{align*}
\sigma_2 e^{-\ii\theta\sg}\sigma_2
&=\sigma_2\begin{pmatrix}e^{-\ii\theta}&0\\0&e^{\ii\theta}\end{pmatrix}\sigma_2\\
&=\begin{pmatrix}e^{\ii\theta}&0\\0&e^{-\ii\theta}\end{pmatrix}
=e^{\ii\theta\sg}.
\end{align*}
因为 $\sigma_2^2=I$, 乘积的约化可在中间插入 $I$:
\begin{align*}
\sigma_2\conj{Y(\conj{z})e^{\ii\theta(\conj{z})\sg}}\sigma_2
&=(\sigma_2\conj{Y(\conj{z})}\sigma_2)
(\sigma_2 e^{-\ii\theta(z)\sg}\sigma_2)\\
&=Y(z)e^{\ii\theta(z)\sg}.
\end{align*}
对变换后的解求导,
\begin{align*}
\partial_x\bigl(\sigma_2\conj{\phi_\pm(\conj{z})}\sigma_2\bigr)
&=\sigma_2\conj{U(\conj{z})}\conj{\phi_\pm(\conj{z})}\sigma_2\\
&=U(z)\sigma_2\conj{\phi_\pm(\conj{z})}\sigma_2.
\end{align*}
这就是同一空间方程; 上式又给出同一端态,
所以可应用 Jost 唯一性.

共轭后的 Jost 解满足同一空间方程和同一归一化, 定理 \ref{thm:jost} 的唯一性给出 \eqref{eq:conjugation}. 在传播方程适用时, $\sigma_2\conj{R}\sigma_2=-R$ 同时保证 $V(z)=\sigma_2\conj{V(\conj{z})}\sigma_2$.


\par\medskip\noindent\textbf{时间矩阵约化中有理项的核对.}\quad 
对介质矩阵逐元计算,
\begin{align*}
\sigma_2\conj{R}\sigma_2
&=\sigma_2
\begin{pmatrix}\eta&-\conj{p}\\-p&-\eta\end{pmatrix}\sigma_2\\
&=\begin{pmatrix}-\eta&p\\\conj{p}&\eta\end{pmatrix}=-R.
\end{align*}
因此有理项经过约化时出现两个负号:
\[
\sigma_2\conj{\left(\dfrac{\ii R}{k(\conj{z})-a}\right)}\sigma_2
=-\dfrac{\ii\sigma_2\conj{R}\sigma_2}{k(z)-a}
=\dfrac{\ii R}{k(z)-a}.
\]
$\dfrac{\ii\Omega}{2}\sg$ 也由
$\conj{\ii}=-\ii$ 与 $\sigma_2\sg\sigma_2=-\sg$
的两个负号保持不变.
对多项式部分, 每个 $V_j$ 都是迹零矩阵,
其对角元纯虚且 $(V_j)_{21}=-\conj{(V_j)_{12}}$.
直接使用一般二阶乘法即有
\[
\sigma_2\conj{V_j}\sigma_2
=\begin{pmatrix}
\conj{(V_j)_{22}}&-\conj{(V_j)_{21}}\\
-\conj{(V_j)_{12}}&\conj{(V_j)_{11}}
\end{pmatrix}=V_j.
\]
例如 $q\conj{q_x}-q_x\conj{q}$ 的共轭是其相反数,
所以 $V_0$ 的对角元也确实纯虚.
这就逐项给出了 $V$ 的共轭约化.

换叶保持 $k$ 而使 $\lambda,\theta$ 变号. 直接由 \eqref{eq:E} 计算,
\begin{equation}\label{eq:Esheet}
Y(h(z))=I+\dfrac{\ii z}{q_0}\sigma_1=Y(z)C(z),\qquad
C(z)\sg=-\sg C(z).
\end{equation}

\par\medskip\noindent\textbf{换叶所需的全部代数恒等式.}\quad 
把 $h(z)$ 代入 $k,\lambda$ 的有理表达式,
\begin{align*}
k(h(z))
&=\dfrac{1}{2}\left(-\dfrac{q_0^2}{z}
-\dfrac{q_0^2}{-\dfrac{q_0^2}{z}}\right)
=\dfrac{1}{2}\left(z-\dfrac{q_0^2}{z}\right)=k(z),\\
\lambda(h(z))
&=\dfrac{1}{2}\left(-\dfrac{q_0^2}{z}
+\dfrac{q_0^2}{-\dfrac{q_0^2}{z}}\right)
=-\dfrac{1}{2}\left(z+\dfrac{q_0^2}{z}\right)=-\lambda(z).
\end{align*}
因为 $F$ 只依赖 $k$, 可得 $F(k(h(z)))=F(k(z))$,
故 $\theta(h(z))=-\theta(z)$.
$Y$ 与 $C$ 的乘法也可直接展开:
\begin{align*}
Y(z)C(z)
&=\left(I-\dfrac{\ii q_0}{z}\sigma_1\right)
\dfrac{\ii z}{q_0}\sigma_1\\
&=\dfrac{\ii z}{q_0}\sigma_1
-\dfrac{\ii q_0}{z}\dfrac{\ii z}{q_0}\sigma_1^2\\
&=I+\dfrac{\ii z}{q_0}\sigma_1=Y(h(z)).
\end{align*}
其逆与两次换叶的一致性分别为
\[
C(z)^{-1}=-\dfrac{\ii q_0}{z}\sigma_1,
\qquad C(z)C(h(z))=I.
\]
由于
\[
\sigma_1\sg=\begin{pmatrix}0&-1\\1&0\end{pmatrix},
\qquad
\sg\sigma_1=\begin{pmatrix}0&1\\-1&0\end{pmatrix},
\]
有 $C\sg=-\sg C$.
把指数写为对角矩阵可逐元确认
\[
e^{\ii\theta\sg}C
=C e^{-\ii\theta\sg},
\qquad
C e^{\ii\theta\sg}=e^{-\ii\theta\sg}C.
\]
因此换叶不是简单地把指数换号后保持列号;
右乘 $C$ 同时交换了两列.

于是 $\phi_\pm(z)C(z)$ 在 $h(z)$ 处满足同一空间方程, 且
\[
Y(z)e^{\ii\theta(z)\sg}C(z)
=Y(h(z))e^{\ii\theta(h(z))\sg}.
\]
再次由唯一性得到 $\phi_\pm$ 的换叶关系. 去除指数时使用
$C(z)e^{\ii\theta(z)\sg}=e^{-\ii\theta(z)\sg}C(z)$,
便得到 $\mu_\pm$ 的关系. 
\par\medskip\noindent\textbf{修正 Jost 关系及四条列关系逐项展开.}\quad 
从 $\phi(h(z))=\phi(z)C(z)$ 出发,
\begin{align*}
\mu(h(z))
&=\phi(h(z))e^{-\ii\theta(h(z))\sg}\\
&=\phi(z)C(z)e^{\ii\theta(z)\sg}\\
&=\phi(z)e^{-\ii\theta(z)\sg}C(z)=\mu(z)C(z).
\end{align*}
共轭关系对 $\mu$ 的证明也保留两个指数负号:
\begin{align*}
\sigma_2\conj{\mu(\conj{z})}\sigma_2
&=\sigma_2\conj{\phi(\conj{z})}
 e^{\ii\theta(z)\sg}\sigma_2\\
&=(\sigma_2\conj{\phi(\conj{z})}\sigma_2)
 (\sigma_2 e^{\ii\theta(z)\sg}\sigma_2)\\
&=\phi(z)e^{-\ii\theta(z)\sg}=\mu(z).
\end{align*}
把右乘 $\sigma_2$ 的结果写成两列,
\[
\conj{\phi(\conj{z})}\sigma_2
=\bigl(\ii\conj{\phi_2(\conj{z})},
-\ii\conj{\phi_1(\conj{z})}\bigr).
\]
再左乘 $\sigma_2$, 得
\[
\phi_1(z)=\ii\sigma_2\conj{\phi_2(\conj{z})},
\qquad
\phi_2(z)=-\ii\sigma_2\conj{\phi_1(\conj{z})}.
\]
右乘 $C$ 则没有两个不同的符号:
\[
(\phi_1,\phi_2)C
=\left(\dfrac{\ii z}{q_0}\phi_2,
\dfrac{\ii z}{q_0}\phi_1\right).
\]
相同乘法适用于 $\mu$.
对解析列比较时, 例如
$\mu_{-,1}(z)$ 与 $\mu_{-,2}(\conj{z})$,
右边的第二列在 $D_-$ 解析, 而共轭后的自变量属于 $D_+$;
所以两边确实在同一个区域上有意义.

右乘矩阵的两列分别给出 \eqref{eq:jost-conjugate-columns}--\eqref{eq:jost-sheet-columns}. 变换后的列具有 \eqref{eq:symmetry-domains} 所确定的解析域, 因此同一 Volterra 唯一性论证在各解析域内成立.
\end{proof}

\begin{proposition}[散射矩阵与反射数据的约化]\label{prop:scattering-symmetry}
散射矩阵满足
\begin{equation}\label{eq:scattering-symmetry}
S(z)=\sigma_2\conj{S(\conj{z})}\sigma_2,\qquad
S(h(z))=\sigma_1S(z)\sigma_1.
\end{equation}
特别地,
\begin{equation}\label{eq:sconj}
s_{22}(z)=\conj{s_{11}(\conj{z})},\qquad
s_{12}(z)=-\conj{s_{21}(\conj{z})},
\end{equation}
\begin{equation}\label{eq:ssheet}
s_{22}(h(z))=s_{11}(z),\qquad s_{12}(h(z))=s_{21}(z),
\end{equation}
\begin{equation}\label{eq:reflection-symmetry}
\rho(z)=-\conj{\widetilde\rho(\conj{z})},\qquad
\rho(h(z))=\widetilde\rho(z).
\end{equation}
透射系数的关系在相应解析域内成立. 在实谱的普通点,
\begin{equation}\label{eq:scattering-unitarity}
S(z)^\dagger S(z)=I,\qquad
|s_{11}(z)|^2+|s_{21}(z)|^2=1.
\end{equation}
\end{proposition}
\begin{proof}

\par\medskip\noindent\textbf{由基底的约化得到散射矩阵的约化.}\quad 
先在 $\conj{z}$ 处写散射关系并取共轭,
\[
\conj{\phi_+(\conj{z})}
=\conj{\phi_-(\conj{z})}\conj{S(\conj{z})}.
\]
左右乘 $\sigma_2$, 并在乘积中间插入
$\sigma_2^2=I$, 得
\begin{align*}
\phi_+(z)
&=(\sigma_2\conj{\phi_-(\conj{z})}\sigma_2)
(\sigma_2\conj{S(\conj{z})}\sigma_2)\\
&=\phi_-(z)(\sigma_2\conj{S(\conj{z})}\sigma_2).
\end{align*}
与 $\phi_+=\phi_-S$ 比较, 左乘可逆的 $\phi_-^{-1}$,
得到共轭约化.
换叶时在 $h(z)$ 处写散射关系,
\[
\phi_+(z)C(z)=\phi_-(z)C(z)S(h(z)).
\]
同时原关系右乘 $C(z)$ 给出
$\phi_+(z)C(z)=\phi_-(z)S(z)C(z)$.
两式消去 $\phi_-$ 后,
\[
C(z)S(h(z))=S(z)C(z),
\qquad
S(h(z))=C(z)^{-1}S(z)C(z).
\]
$C$ 中非零的标量因子在共轭中抵消, 因而只留下
$\sigma_1S\sigma_1$.

将 \eqref{eq:conjugation} 代入 $\phi_+=\phi_-S$, 并利用 $\phi_-$ 的可逆性, 得到 \eqref{eq:scattering-symmetry} 的第一式. 比较
\[
\phi_+(z)C(z)=\phi_-(z)S(z)C(z)
=\phi_-(z)C(z)S(h(z)),
\]
得到
\[
S(h(z))=C(z)^{-1}S(z)C(z)=\sigma_1S(z)\sigma_1.
\]

\par\medskip\noindent\textbf{把两个矩阵约化完全展开为四个矩阵元.}\quad 
共轭约化的右侧实际为
\begin{align*}
\sigma_2\conj{S(\conj{z})}\sigma_2
&=\begin{pmatrix}
\conj{s_{22}(\conj{z})}&-\conj{s_{21}(\conj{z})}\\
-\conj{s_{12}(\conj{z})}&\conj{s_{11}(\conj{z})}
\end{pmatrix}.
\end{align*}
故除了原文列出的两式之外, 另外两式为
\[
s_{11}(z)=\conj{s_{22}(\conj{z})},
\qquad
s_{21}(z)=-\conj{s_{12}(\conj{z})}.
\]
换叶约化则给出
\[
\sigma_1S(z)\sigma_1
=\begin{pmatrix}s_{22}(z)&s_{21}(z)\\s_{12}(z)&s_{11}(z)\end{pmatrix}.
\]
所以四个矩阵元分别是
\begin{align*}
s_{11}(h(z))&=s_{22}(z),&s_{12}(h(z))&=s_{21}(z),\\
s_{21}(h(z))&=s_{12}(z),&s_{22}(h(z))&=s_{11}(z).
\end{align*}
按反射系数的分母逐项相除,
\begin{align*}
\rho(z)
&=\dfrac{-\conj{s_{12}(\conj{z})}}{\conj{s_{22}(\conj{z})}}
=-\conj{\left(\dfrac{s_{12}(\conj{z})}{s_{22}(\conj{z})}\right)}
=-\conj{\widetilde\rho(\conj{z})},\\
\rho(h(z))
&=\dfrac{s_{21}(h(z))}{s_{11}(h(z))}
=\dfrac{s_{12}(z)}{s_{22}(z)}=\widetilde\rho(z).
\end{align*}
负号来自 $\sigma_2$ 的非对角元, 不来自行列式约定的改变.
在非实圆谱上应保留 $\conj{z}$ 这个自变量;
只有 $z$ 为实时才可把它替换为 $z$.

逐元展开即得 \eqref{eq:sconj}--\eqref{eq:ssheet}; 除以相应透射系数得到 \eqref{eq:reflection-symmetry}. 解析透射系数的关系也可直接对 \eqref{eq:wronskian} 应用 Jost 列的约化得到.


\par\medskip\noindent\textbf{实轴上的酉性怎样逐项得到.}\quad 
在实 $z$ 上,
\[
U=\begin{pmatrix}-\ii k&q\\-\conj{q}&\ii k\end{pmatrix},
\qquad
U^\dagger=\begin{pmatrix}\ii k&-q\\\conj{q}&-\ii k\end{pmatrix}=-U.
\]
所以
\begin{align*}
\partial_x(\phi_\pm^\dagger\phi_\pm)
&=(U\phi_\pm)^\dagger\phi_\pm
+\phi_\pm^\dagger U\phi_\pm\\
&=\phi_\pm^\dagger U^\dagger\phi_\pm
+\phi_\pm^\dagger U\phi_\pm=0.
\end{align*}
背景矩阵的乘积为
\begin{align*}
Y^\dagger Y
&=\begin{pmatrix}1&\dfrac{\ii q_0}{z}\\
\dfrac{\ii q_0}{z}&1\end{pmatrix}
\begin{pmatrix}1&-\dfrac{\ii q_0}{z}\\
-\dfrac{\ii q_0}{z}&1\end{pmatrix}\\
&=\begin{pmatrix}1+\dfrac{q_0^2}{z^2}&0\\
0&1+\dfrac{q_0^2}{z^2}\end{pmatrix}=\gamma I.
\end{align*}
由于 $\theta$ 为实数, $e^{\ii\theta\sg}$ 是酉矩阵,
在归一化端
\[
\phi_\pm^\dagger\phi_\pm
\longrightarrow e^{-\ii\theta\sg}\gamma I e^{\ii\theta\sg}=\gamma I.
\]
空间守恒性使该等式对所有有限 $x$ 成立.
代入 $\phi_+=\phi_-S$,
\[
\gamma I=\phi_+^\dagger\phi_+
=S^\dagger\phi_-^\dagger\phi_-S
=\gamma S^\dagger S.
\]
$\gamma=1+\dfrac{q_0^2}{z^2}>0$, 故可除以它.
最后逐元展开得到
\[
S^\dagger S
=\begin{pmatrix}
|s_{11}|^2+|s_{21}|^2&\conj{s_{11}}s_{12}+\conj{s_{21}}s_{22}\\
\conj{s_{12}}s_{11}+\conj{s_{22}}s_{21}&|s_{12}|^2+|s_{22}|^2
\end{pmatrix}=I.
\]
左上元就是原文的标量恒等式, 右上元则给出两列的正交关系.
结合 $\widetilde\rho=-\conj\rho$ 还可复核
\[
|s_{11}|^2(1+|\rho|^2)=1,
\]
这与行列式关系的实轴取值一致.

在实 $z$ 上, $k,\lambda,F$ 为实数, 且 $U^\dagger=-U$. 因而
\[
\partial_x(\phi_\pm^\dagger\phi_\pm)
=\phi_\pm^\dagger(U^\dagger+U)\phi_\pm=0,\qquad
Y^\dagger Y=\gamma I.
\]
振荡归一化给出 $\phi_\pm^\dagger\phi_\pm=\gamma I$. 代入散射关系并除以 $\gamma>0$, 即得 \eqref{eq:scattering-unitarity}.
\end{proof}
